% Generated draft: requires mathematical and source review.
\documentclass[11pt]{article}
\usepackage[margin=25mm]{geometry}
\usepackage{fontspec}
\setmainfont{Latin Modern Roman}
\usepackage{amsmath,amssymb,mathtools,unicode-math}
\setmathfont{Latin Modern Math}
\usepackage{enumitem,needspace,xurl,hyperref}
\setlength{\parindent}{0pt}
\setlength{\parskip}{5pt}
\setlength{\emergencystretch}{3em}
\newcommand{\N}{\mathbb{N}}
\newcommand{\Z}{\mathbb{Z}}
\newcommand{\Q}{\mathbb{Q}}
\newcommand{\R}{\mathbb{R}}
\newcommand{\C}{\mathbb{C}}
\newcommand{\F}{\mathbb{F}}
\newcommand{\A}{\mathbb{A}}
\newcommand{\PP}{\mathbb{P}}
\newcommand{\E}{\mathbb{E}}
\newcommand{\eps}{\varepsilon}
\newcommand{\dd}{\,\mathrm{d}}
\begin{document}
\section*{UnsolvedMath ID: 8800004}
\textbf{Source problem code:} catalogue entry
\subsection{Definitions and mathematical statement}
For a finite field \(\F_q\), a cyclic subgroup \(\langle g\rangle\leq\F_q^\times\) of order \(r\), and \(h\in\langle g\rangle\), the discrete logarithm problem asks for \(x\in\Z/r\Z\) satisfying \(h=g^x\).

Frey's task is to make the \emph{proposed realization} of existing finite-field discrete-logarithm systems inside class groups explicit for practical systems. This is a construction and implementation problem, not just an abstract assertion that two finite groups share a cyclic subgroup. A complete explicit realization would need to specify the relevant class group, effective representations and maps for the original system's data, and how the discrete logarithm relation is represented and recovered. The direction in the statement is from the finite-field system to its class-group realization.

The proposed construction is not included in the excerpt. ``Class group'' could refer to ideal classes of a number field or divisor/Picard classes of a curve, and the source does not identify which. No field family, parameter size, complexity bound, security criterion, or meaning of practical is specified. The supplied summary mentions a pairing-based reduction but gives neither a pairing nor its domain, target, or lifting maps. A group embedding, an isomorphism onto a subgroup, and a pairing reduction are distinct structures; none is silently selected as the requested realization.
\subsection{Short English statement}
Give an explicit practical class-group realization of existing finite-field discrete-logarithm systems. The omitted proposed construction and the intended class groups and efficiency criteria are needed to turn this into a precise algorithmic specification.
\subsection{Sources}
\begin{itemize}
\item Local UnsolvedMath registry, permanent ID 8800004, source code catalogue entry; attributed to Gerhard Frey; no proposal year supplied; update 2026-08-06.
\item IPAM, \emph{Workshop on Number Theory and Cryptography --- Open Problems} (2006), attachment 6219, PDF page 18, Problem 3: \url{https://www.ipam.ucla.edu/programs/workshops/workshop-i-number-theory-and-cryptography-open-problems/?tab=schedule}. Registry access date 2026-07-29; PDF and legacy-PPT visual audit. Provenance reports 19 surviving attachments, comprising 17 PDFs and two legacy PPT decks, with animation repetitions and solved definitions excluded. Current primary-literature status and redistribution terms require review.
\item The supplied partial-progress summary reports an explicit pairing-based reduction but no complete practical self-contained realization in the literature consulted. No additional exact reference or construction is supplied, and that assessment was not independently checked here.
\end{itemize}
\end{document}
